% ============================================================
% Section 4.6 -- Limited Attention to Silence (extension)
% DRAFT inserted 2026-10-08 from playground/notes.md Part IV (plan harmonic-munching-lecun.md).
% Every number is a macro generated by scripts/compute_attention.py and
% scripts/mc_attention_estimator.py; values are provisional until the certificate-date (G2)
% and certified-consumption (G3) re-pulls. Paragraphs tagged [AUTHOR REWRITE] are
% load-bearing and should be rewritten in the author's voice (writing-with-ai.md).
% ============================================================

\subsection{Limited Attention to Silence}
\label{ss:attention}

% [AUTHOR REWRITE] opening
The estimates in Section~\ref{sec:results} leave one result unexplained, and its shape is specific enough to point toward a mechanism. In apartment markets the reform raised the price of a listing that shows no rating, relative to one rated E to G, by more than it raised the price of a verified A-to-D listing. We call a listing that shows no rating \emph{silent}; in the notation of Section~\ref{ss:model_setup} it is a non-disclosing listing, $d=0$. Averaged over the \attNMarketsApt{} apartment markets of Table~\ref{tab:did_results}, the two shifts are \attDthreeApt{} and \attBthreeApt{} log points, and the first exceeds the second in \attNDgtBApt{} of them. In house markets the order is reversed on average, \attDthreeHouse{} against \attBthreeHouse{}. This subsection asks what would have to be true of buyers for silence to be worth more than a verified efficient rating.

One answer is that buyers attend to energy when a listing puts it in front of them, and much less when it does not. We allow for this mechanism without imposing it: full attention remains a special case, and the data decide between the two. The extension changes the buyer side of the model and nothing else. It nests the model of Sections~\ref{ss:model_setup}--\ref{sec:equilibrium} as a special case and preserves its regimes. It leaves $\hat\rho_m$ unchanged as a number but reads it as effective attention to displayed ratings, and it reads the concealment cost off the non-disclosure premium rather than off $\hat\beta_{1,m}\hat e^*_m$ (Proposition~\ref{prop:att_identification}). It also turns the over-identifying wedge $\hat\lambda_m$ of Section~\ref{sec:identification} into an estimate of how much attention silent listings receive.

\subsubsection{Motivation}
\label{ss:att_motivation}

\paragraph{Why the model above cannot reach the apartment wedge.}
The difficulty is not one of calibration. In the model of Section~\ref{ss:model_setup}, a buyer who sees no rating prices the listing at her belief about its quality, and a belief about a binary rating lies between the worst rating and the best. A silent listing therefore sells at most at the price of a verified efficient listing: the ratio of the non-disclosure premium to the energy-efficiency premium is the pool belief, which is bounded by $\bar\theta\le1$ (Lemma~\ref{lem:model_bounds}(i)). The same bound applies to the differences-in-differences whenever the pool belief is stable across the reform.\footnote{The ordering survives exact pricing under the monotonicity condition discussed in Appendix~\ref{app:att_exact}. The levels make the same point: silent apartments sold above verified A-to-D listings by \attLevelGapPreApt{} log points before the reform and by \attLevelGapPostApt{} after, whereas the model prices silence below verified efficiency at any date and for any belief in $[0,1]$. Under Assumption~\ref{assum:no_pre_pricing} the pre-reform gap is non-energy, and the energy component of the post-reform gap is the DiD, \attExcessApt{} log points, still positive.} The apartment wedge of \attLamApt{} log points decomposes accordingly. Of it, \attExcessApt{} points are the amount by which $\hat\Delta$ exceeds $\hat\beta_1$, a region the model cannot enter. Only the remaining \attReachApt{} points, $(1-\hat e^*)\hat\beta_1$, are within reach of any change in how buyers form beliefs about silent listings.

Figure~\ref{fig:att_belief} shows that most apartment markets lie above the ceiling that the core model imposes, whatever belief buyers hold about silent listings.

\begin{figure}[htbp]
    \centering
    \caption{The core model cannot produce the non-disclosure premium}
    \label{fig:att_belief}
    \includegraphics[width=\linewidth]{fig/visuals/fig_att_F2a_belief.pdf}
    \begin{minipage}{0.95\linewidth}
        \vspace{0.5em}
        \footnotesize\textit{Notes.} Each point is a market; both axes are log points. In the core model a silent listing sells at the buyers' belief about its quality times the energy-efficiency premium $\hat\beta_{1,m}$, and a belief about a binary rating lies between zero and one, so the non-disclosure premium $\hat\Delta_m$ cannot exceed $\hat\beta_{1,m}$ (the line). Markets in the shaded region, drawn as filled points, are outside what the core model can produce for any belief about silent listings: \attNDgtBApt{} of \attNMarketsApt{} apartment markets and \attNDgtBHouse{} of \attNMarketsHouse{} house markets. Sampling error in $\hat\beta_{1,m}$ and $\hat\Delta_m$ is not shown.
    \end{minipage}
\end{figure}

\paragraph{Belief-based amendments cannot reach the wedge.}
This rules out the amendments that come to mind first. First, buyers may draw too little inference from silence, as in cursed equilibrium and in experimental evidence on disclosure games \citep[e.g.,][]{eyster_cursed_2005, jin_is_2021, hagenbach_motivated_2025}. Any such rule replaces the Bayesian belief $e^*$ by another number between zero and one, so it stays below the ceiling. It is also quantitatively small here. Sellers without a certificate make up about \attShareUninfApt{} percent of the silent apartment pool, so the Bayesian belief already sits within a few points of the prior, and full cursedness would account for about \attCursedApt{} of the \attLamApt{} points. Second, buyers may over-react to a displayed poor rating. That lowers the price of E-to-G listings, which is the base of both gaps, and it cancels from the comparison between silent and verified efficient listings. Third, the distrust discount on inferred quality allowed in Section~\ref{ss:model_setup} ($\alpha_k<\beta_k$) moves the price of silence down, not up. Observable composition does not resolve the puzzle either, because the indices hold the hedonic reference property fixed within each market and the comparison is a difference-in-differences. What these arguments do not exclude is a change in the \emph{unobserved} composition of the silent pool, most of which consists of sellers without a certificate: if the dwellings that were silent after the reform were of better unobserved quality, the wedge would rise without any change in attention.


% [AUTHOR REWRITE] mechanism paragraphs
What these amendments share is that the buyer facing a silent listing still thinks about its energy performance; she only reasons about it imperfectly. A large body of evidence suggests that consumers often do not think about costs that are not put in front of them, and that when they do not, they act as if those costs were small rather than average. Buyers under-react to sales taxes that are not included in posted prices \citep{chetty_salience_2009}, to road tolls collected electronically \citep{finkelstein_eztax_2009} and to shrouded shipping fees \citep{brown_shrouded_2010}. Energy operating costs are a documented instance in durable-goods markets \citep[e.g.,][]{allcott_consumers_2011, sallee_rational_2014, allcott_evaluating_2015}. Models of focusing, salience and sparsity formalise the same idea, that the attributes a choice set makes salient receive more weight than those it leaves in the background \citep[e.g.,][]{koszegi_model_2013, bordalo_salience_2013, gabaix_sparsity_2014}. In financial markets, limited attention changes how disclosed and undisclosed information is priced \citep[e.g.,][]{hirshleifer_limited_2003, dellavigna_investor_2009}.

Three features of the French setting make this behaviour plausible for apartment buyers in particular. First, a listing that discloses shows the rating as a coloured label and, from January 2022, an estimate of the annual energy bill in euros, so disclosure itself puts energy into the buyer's comparison of dwellings.\footnote{Since 1 January 2022, advertisements for the sale of a dwelling must state the estimated annual energy expenditure for standard use, with the reference year of energy prices, and dwellings rated F or G must carry the mention ``logement \`a consommation \'energ\'etique excessive''. Omitting these mentions exposes the advertiser to fines under Article L126-33 of the Code de la construction et de l'habitation. % [CHECK] cite the implementing decree (D\'ecret n\textsuperscript{o} 2020-1609, 17 Dec.\ 2020) against L\'egifrance before submission.
The obligation also gives a dated test of the mechanism: under limited attention the gap between silent and rated listings should widen again at that date.} A listing that shows nothing does not. Second, heating in many apartment buildings is collective and paid through service charges, so apartment buyers are likely to see energy as part of a monthly charge rather than as an attribute of the dwelling. Third, a missing rating is common: about \attShareUninfApt{} percent of silent apartment listings match no post-reform certificate in the registry. Silence is therefore unremarkable rather than conspicuous.

The mechanism is easiest to see in three steps. Suppose first that buyers attend to energy as often for a silent listing as for a rated one. They then price silence at the expected cost of the silent pool, and the model is exactly the one estimated above. Suppose next that they attend half as often. Half of them still price silence at the expected cost of the pool, while the other half price it at whatever they implicitly assume when energy never comes up. If that default is ``nothing out of the ordinary'', the silent listing gains on every rated listing, efficient ones included. Suppose finally that buyers never attend to energy on a silent listing. It is then priced as if it carried the default cost, which can lie below the cost of an A-rated dwelling, and silence beats every rating.

\subsubsection{Setup}
\label{ss:att_setup}

We embed display-dependent attention in the buyer side of Section~\ref{ss:model_setup}. This lets us read the over-identifying wedge as a measure of the attention silent listings receive. The estimators of Section~\ref{sec:identification} keep their numerical values, but under the extension they measure different objects (Proposition~\ref{prop:att_identification}). Everything not stated below is as in Section~\ref{ss:model_setup}: seller types, the No Falsification assumption, Assumption~\ref{ass:independence}, the payoff~\eqref{eq:payoff} and the demand system~\eqref{eq:demand}.

\paragraph{Environment and timing.}
The environment is that of Section~\ref{ss:model_setup}, with one addition on the buyer side: after seeing the listing, an EPC-sensitive buyer either attends to its energy performance or does not, with a probability that depends on what the listing shows. The timing is as follows.
\begin{enumerate}[label=(\arabic*)]
    \item Nature draws the seller's quality $\theta$ and information status $K$, independently, with $\Pr(\theta=1)=\bar\theta$ and $\Pr(K=1)=\pi$.
    \item The seller lists the property, choosing the display $d$ (only if $K=1$) and the listing price.
    \item A buyer arrives. She is EPC-sensitive with probability $\rho$. If sensitive, she attends to energy with a probability $\omega_d$ that depends on the display.
    \item She values the property given what she attends to, and buys or not.
\end{enumerate}

\paragraph{Payoffs and hiding costs.}
The seller's problem is~\eqref{eq:payoff}, with hiding cost $c(\theta,d)=\mathbf 1[d=0](c_0+\theta c_1)$, interpreted as in Section~\ref{ss:model_setup}. The extension leaves the seller side of the model unchanged. It changes only how buyers price what a listing shows, and therefore what a seller gains by showing nothing.

\paragraph{Buyer valuations and demand.}
Section~\ref{ss:model_setup} writes the valuation of an EPC-sensitive buyer in what we call \emph{bonus form}: she adds $\beta_1$ times her belief that the dwelling is efficient. This form records the difference in energy value between efficient and inefficient dwellings, $\beta_1$, but never the level of the energy cost of either, and it has no place for a buyer who forms no belief at all, because every sensitive buyer enters through some number $\mathbb E[\theta\mid\cdot\,]\in[0,1]$. The ceiling documented in Section~\ref{ss:att_motivation} is thus a property of the representation as much as of the economics. To ask what a buyer pays for a listing whose energy performance she does not consider, we rewrite the valuation in \emph{cost form}. The rewriting changes nothing in the model of Section~\ref{ss:model_setup} (Lemma~\ref{lem:cost_form}); it changes what the model can express. Appendix~\ref{app:attention} explains why the normalisation implicit in the bonus form becomes binding once attention depends on what a listing shows.

A dwelling of quality $\theta$ in market $m$ carries a capitalised energy cost $C_{\theta,m}$, with $C_{0,m}>C_{1,m}\ge0$, measured in the units of the price gaps; we drop the market index when no confusion arises. We take these costs from the construction of the anchor in Assumption~\ref{assum:energy_anchor}, $C_{\theta,m}=\kappa\,p^E a(r,T)\,\bar E_{\theta,m}/P_m$, so that $\beta^{\rm anchor}_{1,m}=C_{0,m}-C_{1,m}$. The extension also uses the cost \emph{level}. Because the conversion factor $\kappa p^Ea(r,T)/P_m$ is common to both qualities, any energy cost below can equivalently be expressed in kWh$_{\rm ep}$/m$^2$/yr. Define the \emph{energy-cost ratio}
\begin{equation}
    \eta_m\;\equiv\;\frac{C_{0,m}}{C_{0,m}-C_{1,m}}\;=\;\frac{\bar E_{0,m}}{\bar E_{0,m}-\bar E_{1,m}}\;>\;1,
    \label{eq:eta}
\end{equation}
which is computed from the registry alone and does not depend on the energy price, the discount rate or the horizon.

\begin{definition}[Gross Value and Cost-Form Valuation]
\label{def:cost_form}
The \emph{gross value} of a dwelling with attributes $X$ is
\begin{equation}
    \bar V(X)\;\equiv\;\beta_0+\sum_n\beta_nX_n+C_0,
    \label{eq:gross_value}
\end{equation}
its value before any energy cost is subtracted. A buyer's \emph{cost-form valuation} of a listing is
\begin{equation}
    u\;=\;\bar V(X)-\mathcal C,
    \label{eq:cost_valuation}
\end{equation}
where $\mathcal C$ is the energy cost she subtracts, given what the listing shows and what she attends to.
\end{definition}

The gross value adds back $C_0$ because, in the bonus form, $\beta_0+\sum_n\beta_nX_n$ is not the value of a dwelling without an energy bill. It is the value of an inefficient dwelling, whose energy cost is already netted out of the intercept. Adding $C_0$ moves the zero point of the valuation from the worst dwelling to a dwelling with no energy cost, so that every dwelling, inefficient ones included, carries its energy cost explicitly. The level of $\bar V(X)$ never matters for a buyer who treats every listing alike, because it cancels from every price gap; what the cost form adds is the ability to compare what a buyer subtracts on one listing with what she subtracts on another.

\begin{lemma}[Bonus and Cost Forms]
\label{lem:cost_form}
The cost-form valuation~\eqref{eq:cost_valuation} reproduces the valuations of Section~\ref{ss:model_setup} exactly, with $\beta_1=C_0-C_1$. An EPC-sensitive buyer subtracts $\mathcal C=C_\theta$ for a disclosed rating of quality $\theta$ and $\mathcal C=\bar C^{\rm sil}$ for a silent listing, where
\begin{equation}
    \bar C^{\rm sil}\;\equiv\;\mathbb E[C\mid d=0]\;=\;e\,C_1+(1-e)\,C_0
    \label{eq:csil}
\end{equation}
is the expected energy cost of the silent pool under the belief $e$ of~\eqref{eq:belief_general}. An EPC-insensitive buyer subtracts $\mathcal C=C^{\rm ins}\equiv\bar V(X)-v_0(X)$, the same amount for every listing with attributes $X$. Moreover, $C_0$ is the only constant which, added to $\beta_0+\sum_n\beta_nX_n$, makes the two forms agree for both qualities.
\end{lemma}

\noindent\textit{Proof.} See Appendix~\ref{app:attention}. $\Box$

\begin{definition}[Default Energy Cost]
\label{def:default_cost}
The \emph{default energy cost} $C^{\rm def}\ge0$ is the energy cost an EPC-sensitive buyer subtracts from a listing's value when she does not attend to energy. It is measured in the same units as $C_\theta$.
\end{definition}

The term follows the limited-attention literature, in which an agent who does not attend to an attribute replaces it with a default value rather than with an inference \citep{gabaix_sparsity_2014, gabaix_behavioral_2019}. The default is therefore not a belief about the dwelling's rating. It is what the buyer implicitly assumes when the question of energy never arises. Three benchmark values fix ideas. With $C^{\rm def}=0$ no energy cost is priced at all, which does not mean the buyer believes the dwelling consumes nothing, only that energy plays no part in her valuation. With $C^{\rm def}$ equal to the prior mean cost $\bar\theta C_1+(1-\bar\theta)C_0$, the inattentive buyer behaves like one who assumes an average dwelling, which is cursedness by another name. With $C^{\rm def}=\bar C^{\rm sil}$ she happens to subtract exactly what an attentive buyer would infer, and attention becomes irrelevant. The same default applies to rated and silent listings, because a buyer who does not attend to energy, by definition, does not use what the listing shows about it. Which value is relevant is an empirical question: Proposition~\ref{prop:att_incentives}(iii) shows that the extension can explain the apartment wedge only if $C^{\rm def}<C_1$, and Proposition~\ref{prop:att_identification}(iii) shows what the data can say about its level.

\begin{assumption}[Display-Dependent Attention]
\label{assum:attention}
An EPC-sensitive buyer attends to energy with probability $\omega_1$ when the listing displays a rating and with probability $\omega_0\le\omega_1$ when it does not. An attentive buyer subtracts the expected energy cost given what the listing shows: $C_\theta$ if a rating is displayed, $\bar C^{\rm sil}$ otherwise. An inattentive buyer subtracts the default energy cost $C^{\rm def}$ of Definition~\ref{def:default_cost}, whatever the listing shows.
\end{assumption}

We call $\psi\equiv\omega_0/\omega_1\in[0,1]$ the \emph{relative attention to silence}: the attention a listing without a rating receives, per unit of the attention a rated listing receives. We call $\rho\,\omega_d$ the \emph{effective attention} to listings with display $d$, the share of all buyers who price energy when they look at such a listing. Table~\ref{tab:att_buyers} records what each buyer subtracts.

\begin{table}[htbp]
\centering
\caption{Energy Cost Subtracted, by Buyer Type and Listing Display}
\label{tab:att_buyers}
\small\renewcommand{\arraystretch}{1.5}
\begin{tabular*}{\textwidth}{@{\extracolsep{\fill}}lccc@{}}
\toprule\toprule
\textbf{Buyer} & \textbf{Probability} & \textbf{Displayed rating $\theta$} & \textbf{Silent listing} \\
\midrule
EPC-insensitive & $1-\rho$ & $C^{\rm ins}$ & $C^{\rm ins}$ \\
EPC-sensitive, attentive & $\rho\,\omega_d$ & $C_\theta$ & $\bar C^{\rm sil}$ \\
EPC-sensitive, inattentive & $\rho\,(1-\omega_d)$ & $C^{\rm def}$ & $C^{\rm def}$ \\
\bottomrule\bottomrule
\end{tabular*}

\vspace{0.5em}
\raggedright
\footnotesize
\textit{Notes:} $\omega_d=\omega_1$ for a displayed rating and $\omega_0$ for a silent listing. $C^{\rm ins}\equiv\bar V(X)-v_0(X)$ (Lemma~\ref{lem:cost_form}) does not depend on the rating or on whether one is displayed, so it cancels from every price gap.
\end{table}

Each buyer values a listing by~\eqref{eq:cost_valuation}, with $\mathcal C$ given by Table~\ref{tab:att_buyers}. The transaction probability generalises~\eqref{eq:demand} to three buyer segments,
\begin{equation}
    \Gamma(p^{\rm list}\mid d)=\rho\Big[\omega_d\,G\big(\bar V-\mathcal C_d-p^{\rm list}\big)+(1-\omega_d)\,G\big(\bar V-C^{\rm def}-p^{\rm list}\big)\Big]+(1-\rho)\,G\big(v_0-p^{\rm list}\big),
    \label{eq:demand_att}
\end{equation}
where $\mathcal C_d$ is $C_\theta$ for a displayed rating and $\bar C^{\rm sil}$ for a silent listing. All results below are stated under the linear pricing approximation used to derive~\eqref{eq:premium_derived}, in which the equilibrium price equals the buyer-share-weighted mean valuation. Whether their orderings survive the exact demand~\eqref{eq:demand_att} is discussed, not proven, in Appendix~\ref{app:attention}.

Write $p_{\rm H}$, $p_{\rm L}$ and $p_{\rm sil}$ for the equilibrium log prices, at the reference property $X$, of a listing displaying an efficient rating (A to D), a listing displaying an inefficient rating (E to G), and a silent listing. In the notation of Section~\ref{sec:equilibrium}, $p_{\rm sil}-p_{\rm L}$ is the premium $p_0^*-p_1^*$.

\begin{lemma}[Prices and the Perceived Cost of Silence]
\label{lem:att_prices}
Under Assumption~\ref{assum:attention}, the energy-efficiency premium $b$ and the non-disclosure premium $\Delta$ are
\begin{align}
    b&\;\equiv\;p_{\rm H}-p_{\rm L}\;=\;\rho\,\omega_1\,\beta_1, \label{eq:b_att}\\
    \Delta&\;\equiv\;p_{\rm sil}-p_{\rm L}\;=\;\rho\,\omega_1\,(C_0-C^*),\qquad C^*\;\equiv\;C^{\rm def}+\psi\,\big(\bar C^{\rm sil}-C^{\rm def}\big). \label{eq:premium_att}
\end{align}
More generally, a listing displaying a rating with energy cost $C$ sells at $p_{\rm sil}-\rho\,\omega_1\,(C-C^*)$.
\end{lemma}

\noindent\textit{Proof.} See Appendix~\ref{app:attention}. $\Box$

$C^*$ is the energy cost the market behaves as if a silent listing carries: a weighted average of the default and the rational cost of the silent pool, with weight $\psi$ on the latter. The default cancels from the energy-efficiency premium because both listings in that comparison display a rating; it does not cancel from the non-disclosure premium, because a silent listing draws attention $\omega_0$ while a rated listing draws $\omega_1$. We keep $\beta_1=C_0-C_1$ for the valuation coefficient of Section~\ref{ss:model_setup} and write $b$ for the premium. With $\omega_1=1$, $b=\rho\beta_1$, the premium of Section~\ref{sec:equilibrium}, whose reform-induced change $\hat\beta_{1,m}$ estimates.

\paragraph{Beliefs and the effective belief.}
Of the three buyer types in Table~\ref{tab:att_buyers}, only the attentive EPC-sensitive buyer forms a belief about a silent listing. Her belief is the Bayesian $e(\sigma_0,\sigma_1)$ of~\eqref{eq:belief_general}, unchanged: the share of efficient dwellings among silent listings. The insensitive buyer does not care about energy, and the inattentive sensitive buyer would care but does not think about it on this listing, so she subtracts the default instead. The Bayesian belief therefore reaches prices only through attentive buyers.

In Section~\ref{sec:equilibrium} everything the market thinks about silence is summarised by one number, because the ratio of the non-disclosure premium to the energy-efficiency premium \emph{is} the belief. Under limited attention prices aggregate a belief held by some buyers and a default used by others, and no single buyer's belief summarises them. We define the number that plays the role $e$ played before.

\begin{definition}[Effective Belief]
\label{def:effective_belief}
The \emph{effective belief} is the ratio of the non-disclosure premium to the energy-efficiency premium, $\tilde e\equiv\Delta/b$, so that $\Delta=b\,\tilde e$ as $\Delta=\rho\beta_1 e$ in Section~\ref{sec:equilibrium}.
\end{definition}

Equivalently, $\tilde e$ is the belief that rationalises observed prices within the model of Section~\ref{sec:equilibrium}: if every EPC-sensitive buyer attended to energy and held belief $\tilde e$ about every silent listing, the non-disclosure premium would be exactly $\Delta$. That model confines beliefs to $[0,1]$. An effective belief above one is therefore a direct statement that no belief, Bayesian or behavioural, can account for the premium in that model, and the gap between $\tilde e$ and the registry-measured Bayesian belief $e^*$ is precisely what limited attention has to explain.

The effective belief follows in four steps. By Lemma~\ref{lem:att_prices}, the factor $\rho\omega_1$ cancels from the ratio, so $\tilde e=(C_0-C^*)/(C_0-C_1)$. This places the perceived cost of silence on a \emph{quality scale}, $q(C)\equiv(C_0-C)/(C_0-C_1)$, on which an inefficient dwelling's cost maps to 0 and an efficient one's to 1. Because $q$ is linear and $C^*=\psi\,\bar C^{\rm sil}+(1-\psi)\,C^{\rm def}$, the effective belief is the same weighted average of the two points, $\tilde e=\psi\,q(\bar C^{\rm sil})+(1-\psi)\,q(C^{\rm def})$. By~\eqref{eq:csil}, $q(\bar C^{\rm sil})=e$: the cost an attentive buyer subtracts maps back exactly to her Bayesian belief. The point at which an inattentive buyer places every listing is the \emph{default-implied belief}
\begin{equation*}
    e^{\rm def}\;\equiv\;q\big(C^{\rm def}\big)\;=\;\frac{C_0-C^{\rm def}}{C_0-C_1}\;=\;\eta-\frac{C^{\rm def}}{C_0-C_1},
\end{equation*}
which equals $\eta$ when $C^{\rm def}=0$. Combining the steps,
\begin{equation}
    \tilde e\;=\;\psi\,e+(1-\psi)\,e^{\rm def}.
    \label{eq:etilde}
\end{equation}

The effective belief is a weighted average of the belief of buyers who think about silence and the default of those who do not, with weight $\psi$ on the belief. The two pieces live on different ranges. The Bayesian belief is a probability, so $e\le1$. The default-implied belief is not: it exceeds one whenever $C^{\rm def}<C_1$, that is, whenever a buyer who ignores energy implicitly assumes a lower cost than that of an average efficient dwelling. A buyer who prices no energy cost at all treats a silent listing as better than a verified efficient one, because even a verified efficient listing shows a positive cost. With $\psi=1$ the model is that of Section~\ref{sec:equilibrium}; with $\psi=0$ the composition of the silent pool plays no role in prices; and $\tilde e>1$, hence $\Delta>b$, is possible as soon as $\psi<1$ and $e^{\rm def}>1$.

\begin{lemma}[Effective-Belief Monotonicity]
\label{lem:etilde_monotone}
$\tilde e(\sigma_0,\sigma_1)$ is weakly decreasing in $\sigma_0$ and weakly increasing in $\sigma_1$, strictly so if and only if $\psi>0$, with $\partial\tilde e/\partial\sigma_j=\psi\,\partial e/\partial\sigma_j$.
\end{lemma}

\noindent\textit{Proof.} In~\eqref{eq:etilde} only $e$ depends on $(\sigma_0,\sigma_1)$. Differentiate and apply Lemma~\ref{lem:belief_monotone}. $\Box$

In Section~\ref{sec:equilibrium} the hiding rate is pinned down by a feedback loop: more concealment worsens the pool, the worse pool lowers the belief, and the lower belief lowers the gain from further concealment. Lemma~\ref{lem:etilde_monotone} says that under limited attention this loop runs at a fraction $\psi$ of its full strength. Inattentive buyers keep subtracting the default whatever the pool contains.

\subsubsection{Equilibrium}
\label{ss:att_equilibrium}

\paragraph{Disclosure incentives.}
An informed seller compares the price of silence with the price of her own displayed rating. By Lemma~\ref{lem:att_prices}, the gross gain from hiding for an informed seller of quality $\theta$ is
\begin{equation}
    g_0\;=\;b\,\tilde e,\qquad g_1\;=\;b\,(\tilde e-1),\qquad g_0-g_1=b.
    \label{eq:gains}
\end{equation}
The low type's gain is the non-disclosure premium itself, $g_0=\Delta$. Section~\ref{sec:equilibrium} writes this gain as $\Delta(e)=\rho\beta_1e$; under limited attention it becomes
\begin{equation}
    \Delta_\psi(e)\;\equiv\;b\,\big[\psi\,e+(1-\psi)\,e^{\rm def}\big]\;=\;\omega_1\Big[\psi\,\Delta(e)+(1-\psi)\,\rho\beta_1\,e^{\rm def}\Big],
    \label{eq:Delta_psi}
\end{equation}
which equals $\Delta(e)$ when $\omega_1=\psi=1$. This one substitution carries the equilibrium analysis of Section~\ref{sec:equilibrium} over to limited attention: the hiding gain is a weighted average of the gain in Section~\ref{sec:equilibrium}, which depends on the pool, and a gain against the default, which does not.

\begin{proposition}[Disclosure Incentives under Limited Attention]
\label{prop:att_incentives}
Under Assumption~\ref{assum:attention}:
\begin{enumerate}[label=(\roman*)]
    \item an informed seller with $\theta=0$ strictly prefers non-disclosure if and only if $\tilde e>c_0/b$;
    \item an informed seller with $\theta=1$ strictly prefers non-disclosure if and only if $\tilde e>1+(c_0+c_1)/b$, which requires that silence sell above a verified efficient rating;
    \item if $C^{\rm def}<\bar C^{\rm sil}$, silence sells above a verified efficient rating ($\tilde e>1$, equivalently $C^*<C_1$) if and only if $\psi<(C_1-C^{\rm def})/(\bar C^{\rm sil}-C^{\rm def})$, which requires $C^{\rm def}<C_1$;
    \item if $\psi=1$, then $\tilde e=e$ for any $\omega_1$ and $C^{\rm def}$; conditions (i)--(ii) are then the net-of-cost comparisons underlying Proposition~\ref{prop:equilibrium}, and with $c_0=c_1=0$ they reduce to Proposition~\ref{prop:premium}.
\end{enumerate}
\end{proposition}

\noindent\textit{Proof.} See Appendix~\ref{app:attention}. $\Box$

Inattention that is the same for rated and silent listings is equivalent to having fewer EPC-sensitive buyers: it scales both premia together and leaves the model of Section~\ref{sec:equilibrium} intact (part iv). Differential attention matters only if the default is optimistic (part iii); a buyer who, without thinking about energy, implicitly assumes an average dwelling is a cursed buyer by another name and stays below the ceiling.\footnote{If the default equals the prior mean cost $\bar\theta C_1+(1-\bar\theta)C_0$, it is at least $C_1$, and silence never outprices a verified efficient rating. If it equals the rational cost of the silent pool, attention is irrelevant. The evidence on shrouded costs cited above finds under-weighting toward zero rather than toward the mean, which supports an optimistic default. Housing-specific evidence on this point does not yet exist, and it is the assumption on which the extension most depends.} And the high-quality seller loses her automatic preference for disclosure once silence outprices a verified efficient rating, so that only the hiding cost stops her from hiding (part ii). Applied to the seven letters instead of the two quality levels, Lemma~\ref{lem:att_prices} also predicts where silence sits among ratings: the gap between a rated listing and a silent one should be linear in the letter's energy cost and cross zero at $C^*$. The model of Section~\ref{sec:equilibrium} places that crossing between D and E; limited attention with an optimistic default places it at or above A.

\begin{corollary}[High-Quality Sellers Disclose outside Complete Adverse Selection]
\label{cor:att_hq}
In any equilibrium with $\sigma_0^*<1$, $\sigma_1^*=0$. In Complete Adverse Selection ($\sigma_0^*=1$), high-quality hiding is possible only if $\tilde e(1,\sigma_1)>1+(c_0+c_1)/b$ for some $\sigma_1$, which requires that silence sell above a verified efficient rating.
\end{corollary}

\noindent\textit{Proof.} See Appendix~\ref{app:attention}. $\Box$

Corollary~\ref{cor:hq_disclose} survives, but for a different reason. In Section~\ref{sec:equilibrium} high-quality sellers disclose because silence costs them money. Here silence may pay them too, but by~\eqref{eq:gains} it always pays them $b$ less than it pays low-quality sellers, so with a common hiding cost the low type is the first to want to hide. Whenever $\sigma_1^*=0$, the belief reduces to~\eqref{eq:belief_simplified} and the effective belief to $\tilde e(\sigma_0)=\psi\,e(\sigma_0)+(1-\psi)\,e^{\rm def}$.

\paragraph{Equilibrium characterization.}
By analogy with~\eqref{eq:Phi}, define the net hiding gain under limited attention,
\begin{equation}
    \Phi_\psi(\sigma_0)\;\equiv\;\Delta_\psi\big(e(\sigma_0)\big)-c_0\;=\;b\Big[\psi\,\frac{(1-\pi)\bar\theta}{\pi(1-\bar\theta)\sigma_0+(1-\pi)}+(1-\psi)\,e^{\rm def}\Big]-c_0,
    \label{eq:Phi_psi}
\end{equation}
and write $e(1)\equiv(1-\pi)\bar\theta/(1-\pi\bar\theta)$ for the belief when every informed low-quality seller hides.

\begin{proposition}[Equilibrium under Limited Attention]
\label{prop:att_equilibrium}
Under Assumption~\ref{assum:attention}, an equilibrium exists, and it is characterised as follows.
\begin{enumerate}[label=(\roman*)]
    \item If $c_0>\Delta_\psi(\bar\theta)$, the unique equilibrium is Cost-Deterred Hiding, with $\sigma_0^*=0$ and $e^*=\bar\theta$.
    \item If $c_0<\Delta_\psi(e(1))$, every informed low-quality seller hides, $\sigma_0^*=1$ (Complete Adverse Selection). Let $\bar g_1(\sigma_1)\equiv b\,[\tilde e(1,\sigma_1)-1]$ be the high type's gain. Then $\sigma_1^*=0$, with $e^*=e(1)$, is an equilibrium if and only if $\bar g_1(0)\le c_0+c_1$, and $\sigma_1^*=1$ is an equilibrium if and only if $\bar g_1(1)\ge c_0+c_1$. Both can hold, so equilibrium need not be unique in this regime. If $\psi=1$, only $\sigma_1^*=0$ survives.
    \item If $\Delta_\psi(e(1))\le c_0\le\Delta_\psi(\bar\theta)$, the equilibrium is Partial Pooling. If $\psi>0$ it is unique, with
    \begin{equation}
        e^*\;=\;\frac{c_0/b-(1-\psi)\,e^{\rm def}}{\psi},\qquad \sigma_0^*\;=\;\frac{1-\pi}{\pi(1-\bar\theta)}\Big(\frac{\bar\theta}{e^*}-1\Big).
        \label{eq:att_pp}
    \end{equation}
    If $\psi=0$, Partial Pooling requires $c_0=b\,e^{\rm def}$ exactly, and then every $\sigma_0\in[0,1]$ is an equilibrium.
\end{enumerate}
\end{proposition}

\noindent\textit{Proof.} See Appendix~\ref{app:attention}. $\Box$

\begin{table}[htbp]
\centering
\caption{Equilibrium Regimes under Limited Attention}
\label{tab:att_equilibria}
\small\renewcommand{\arraystretch}{1.5}
\begin{tabular*}{\textwidth}{@{\extracolsep{\fill}}lccc@{}}
\toprule\toprule
\textbf{Equilibrium} & \textbf{Condition} & \textbf{Strategy $\sigma_0^*$} & \textbf{Non-disclosure premium $\Delta^*$} \\
\midrule
Cost-Deterred Hiding & $c_0>\Delta_\psi(\bar\theta)$ & $0$ & $b\,[\psi\bar\theta+(1-\psi)e^{\rm def}]$ \\
Partial Pooling & $\Delta_\psi(e(1))\le c_0\le\Delta_\psi(\bar\theta)$ & \eqref{eq:att_pp} & $c_0$ \\
Complete Adverse Selection & $c_0<\Delta_\psi(e(1))$ & $1$ & $b\,[\psi e(1)+(1-\psi)e^{\rm def}]$ \\
\bottomrule\bottomrule
\end{tabular*}

\vspace{0.5em}
\raggedright
\footnotesize
\textit{Notes:} The premium in Complete Adverse Selection is evaluated at $\sigma_1^*=0$. At $\psi=1$ and $\omega_1=1$ the table reduces to Table~\ref{tab:equilibria}.
\end{table}

% [AUTHOR REWRITE] lead-in to the proposition on seller types
The three regimes keep their logic, and the central result of Section~\ref{sec:equilibrium} survives intact: in Partial Pooling the non-disclosure premium equals the hiding cost, because the marginal hider is indifferent. That statement concerns the marginal seller, not buyers, so it holds whatever buyers attend to. What limited attention changes is where silence sits among the ratings, and with it which seller gains most from silence. The next proposition derives the price of each type of listing and the payoff of each type of seller. Uninformed sellers cannot hide a rating they do not have, so they bear no hiding cost.

\begin{proposition}[Prices and Payoffs by Seller Type]
\label{prop:att_seller_types}
Under Assumption~\ref{assum:attention}, with $b=\rho\omega_1\beta_1>0$, measure prices and payoffs relative to a listing displaying an inefficient rating, and consider an equilibrium with $\sigma_1^*=0$. A listing displaying an efficient rating sells at $b$, and a silent listing at
\begin{equation}
    \Delta^*\;=\;b\,\tilde e^*\;=\;\min\Big\{\max\big\{c_0,\,\Delta_\psi(e(1))\big\},\;\Delta_\psi(\bar\theta)\Big\},
    \label{eq:att_premium_trunc}
\end{equation}
the hiding cost truncated to the Partial Pooling band. An informed high-quality seller earns $b$, an informed low-quality seller earns $\max\{\Delta^*-c_0,\,0\}$, and an uninformed seller of either quality earns $\Delta^*$ (Table~\ref{tab:att_seller_types}). By regime:
\begin{itemize}
    \item in \textbf{Cost-Deterred Hiding}, $\Delta^*=\Delta_\psi(\bar\theta)<c_0$, and every informed low-quality seller displays her rating and earns zero;
    \item in \textbf{Partial Pooling}, $\Delta^*=c_0$ and $\tilde e^*=c_0/b$. Every informed low-quality seller earns zero, whether she hides or displays, and the uninformed seller earns $c_0$. A silent listing sells above a verified efficient rating if and only if $c_0>b$. Payoffs are then ordered uninformed, informed high quality, informed low quality; if $0<c_0<b$ they are ordered informed high quality, uninformed, informed low quality;
    \item in \textbf{Complete Adverse Selection}, $\Delta^*=\Delta_\psi(e(1))>c_0$, and every informed low-quality seller hides and earns $\Delta^*-c_0>0$, which is $c_0$ less than the uninformed seller.
\end{itemize}
In every regime, a silent listing sells above a verified efficient rating, $\Delta^*>b$, if and only if $e^{\rm def}>1$ and
\begin{equation}
    \psi\;<\;\bar\psi(e^*)\;\equiv\;\frac{e^{\rm def}-1}{e^{\rm def}-e^*},
    \label{eq:att_psibar}
\end{equation}
the threshold of Proposition~\ref{prop:att_incentives}(iii) evaluated at the equilibrium belief. The uninformed seller then earns more than an informed high-quality seller and at least as much as any informed seller. In an equilibrium with $\sigma_1^*>0$, which exists only in Complete Adverse Selection, $\Delta^*\ge b+c_0+c_1$.
\end{proposition}

\noindent\textit{Proof.} See Appendix~\ref{app:attention}. $\Box$

\begin{table}[htbp]
\centering
\caption{Prices and Payoffs by Listing and Seller Type under Limited Attention}
\label{tab:att_seller_types}
\small\renewcommand{\arraystretch}{1.5}
\begin{tabular*}{\textwidth}{@{\extracolsep{\fill}}lccc@{}}
\toprule\toprule
 & \textbf{Cost-Deterred Hiding} & \textbf{Partial Pooling} & \textbf{Complete Adverse Selection} \\
\midrule
\multicolumn{4}{l}{\textit{A. Price, relative to a displayed inefficient rating}} \\
Displayed inefficient rating (E to G) & $0$ & $0$ & $0$ \\
Displayed efficient rating (A to D) & $b$ & $b$ & $b$ \\
Silent listing, $\Delta^*$ & $\Delta_\psi(\bar\theta)$ & $c_0$ & $\Delta_\psi(e(1))$ \\
\midrule
\multicolumn{4}{l}{\textit{B. Net payoff, relative to a displayed inefficient rating}} \\
Informed, high quality & $b$ & $b$ & $b$ \\
Informed, low quality, displays & $0$ & $0$ & -- \\
Informed, low quality, hides & -- & $0$ & $\Delta_\psi(e(1))-c_0$ \\
Uninformed, either quality & $\Delta_\psi(\bar\theta)$ & $c_0$ & $\Delta_\psi(e(1))$ \\
\midrule
Silence above an efficient rating iff & $\psi<\bar\psi(\bar\theta)$ & $c_0>b$ & $\psi<\bar\psi(e(1))$ \\
\bottomrule\bottomrule
\end{tabular*}

\vspace{0.5em}
\raggedright
\footnotesize
\textit{Notes:} Equilibria with $\sigma_1^*=0$. A dash marks a seller type that is not on the equilibrium path; in Cost-Deterred Hiding a low-quality seller who deviated to silence would earn $\Delta_\psi(\bar\theta)-c_0<0$. Displayed prices do not depend on beliefs (Lemma~\ref{lem:att_prices}) and are defined in every regime. $\bar\psi(e)$ is defined in~\eqref{eq:att_psibar} for $e^{\rm def}>1$; if $e^{\rm def}\le1$ silence never sells above an efficient rating. Uninformed sellers bear no hiding cost; relative to her own verified price, an uninformed high-quality seller earns $\Delta^*-b$. In Complete Adverse Selection with $\sigma_1^*>0$ (Proposition~\ref{prop:att_equilibrium}), $\Delta^*\ge b+c_0+c_1$ and an informed high-quality seller hides and earns $\Delta^*-c_0-c_1\ge b$. At $\psi=\omega_1=1$ the silent price is $\rho\beta_1e^*<b$ in every column (Table~\ref{tab:equilibria}).
\end{table}

\begin{corollary}[Where Silence Sits among Ratings]
\label{cor:att_ordering}
Suppose $c_0+c_1>0$. If $e^{\rm def}>1$ and $\psi<\bar\psi(e(1))$, a silent listing sells above a verified efficient rating in every equilibrium, whatever the hiding cost. If $e^{\rm def}\le1$ or $\psi\ge\bar\psi(\bar\theta)$, it sells at most at the price of a verified efficient rating in every equilibrium, whatever the hiding cost. Otherwise, in any equilibrium with $\sigma_1^*=0$, it sells above a verified efficient rating if and only if $c_0>b$. In the model of Section~\ref{sec:equilibrium}, $\psi=1$ and $0<\Delta^*<b$ in every equilibrium.
\end{corollary}

\noindent\textit{Proof.} See Appendix~\ref{app:attention}. $\Box$

% [AUTHOR REWRITE] empirical reading of the proposition
The proposition gives the two segments one reading. In apartments the non-disclosure premium, \attDthreeApt{} log points, is about twice the energy-efficiency premium, \attBthreeApt{}: the effective belief is \attRatioApt{}, and silent listings sell above verified efficient ones, which sell above inefficient ones. In houses the effective belief is \attRatioHouse{} and silence sells between the two ratings. These orderings are read directly from the two estimated premia. Which case of Corollary~\ref{cor:att_ordering} each segment falls in also depends on the normalisation of the default ($C^{\rm def}=0$): at the estimated attention apartments are in the first case ($\hat\psi_\delta=\attPsiApt{}$ against a threshold of \attPsiBarOneApt{}), so silence outprices a verified efficient rating whatever the hiding cost, and houses are in the second ($\attPsiHouse{}$ against \attPsiBarThetaHouse{}). The model with one $\psi$ per segment therefore predicts the first ordering in every apartment market and the second in every house market. The market-level estimates are noisy and $\psi$ is common within a segment, so the counts differ: silence sells above a verified efficient rating in \attNDgtBApt{} of \attNMarketsApt{} apartment markets and \attNDgtBHouse{} of \attNMarketsHouse{} house markets, between the two displayed ratings in \attNBetweenApt{} and \attNBetweenHouse{}, and below a displayed inefficient rating in \attRawNNegApt{} and \attRawNNegHouse{}. The last cannot arise in Partial Pooling and reflects sampling error or heterogeneity that we do not model. The payoff ranking among informed sellers is stated at the hiding cost that reproduces the observed hiding rate, which at the aggregate and in Partial Pooling is the premium itself, $c_0=\hat\Delta$ (Proposition~\ref{prop:att_identification}(iv)); the price ordering needs only $\hat\Delta$ against $\hat\beta_1$. Read through Partial Pooling, an uninformed apartment seller earns \attDthreeApt{} log points over a displayed inefficient rating, an informed high-quality seller \attBthreeApt{}, and an informed low-quality seller nothing, whether she hides or not. Silence paying more than a verified rating does not make high-quality sellers hide, which is an implication of Partial Pooling and not a separate finding: their gain, \attExcessApt{} log points, falls short of the hiding cost \attDthreeApt{} that keeps low-quality sellers indifferent. The uninformed seller's rent is a property of the maintained assumptions, that uninformed sellers have no strategy, bear no hiding cost and are exogenous in number, and not a finding.

Limited attention also changes the location and the width of the band of hiding costs that sustains Partial Pooling. Both bounds shift up by $(1-\psi)\,b\,e^{\rm def}$, the part of the gain that does not depend on the pool, and the width shrinks by the factor $\psi$.

\begin{corollary}[The Partial Pooling Band]
\label{cor:att_band}
The interval of hiding costs that supports Partial Pooling is $[\Delta_\psi(e(1)),\,\Delta_\psi(\bar\theta)]$, of width $\psi\,b\,[\bar\theta-e(1)]$. Its width goes to zero as $\psi\to0$.
\end{corollary}

\noindent\textit{Proof.} Direct from~\eqref{eq:Delta_psi} and Proposition~\ref{prop:att_equilibrium}, as in Corollary~\ref{cor:rho}(ii). $\Box$

The corollary has a sharp empirical counterpart. At the apartment estimates of Section~\ref{ss:att_estimates}, taking the reform-period estimate $\hat\psi_\delta$ as a proxy for $\psi$ (exact under the maintained assumption of Proposition~\ref{prop:att_identification}), the band spans \attBandLoApt{} to \attBandHiApt{} log points, a width of \attBandWidthApt{}, against \attBandWidthBaseApt{} in the model of Section~\ref{sec:equilibrium}. The mean premium lies inside the aggregate band by construction of the aggregation, which is not evidence. Market by market, with a common $\psi$, the premium falls inside the market's own band in \attShareInBandApt{} percent of apartment markets: a single hiding cost per market would have to fall inside a window of about \attBandWidthApt{} log points for an apartment market to show interior hiding, yet the registry shows interior hiding rates in essentially every market. In houses the band is wider, \attBandWidthHouse{} log points, but only \attShareInBandHouse{} percent of house markets fall inside their own band. A single hiding cost per market therefore generates interior hiding only within a narrow window of costs, particularly in apartments; we do not model heterogeneity in hiding costs across sellers, and we read the band as a limitation of the single-cost specification.

\paragraph{Comparative statics.}
The next proposition collects how the equilibrium responds to attention, to the default, and to the parameters of Section~\ref{sec:equilibrium}.

\begin{proposition}[Comparative Statics under Limited Attention]
\label{prop:att_cs}
Suppose $C^{\rm def}<\bar C^{\rm sil}$.
\begin{enumerate}[label=(\roman*)]
    \item In Partial Pooling, $\partial\sigma_0^*/\partial\psi<0$, $\partial\sigma_0^*/\partial C^{\rm def}<0$, $\partial\sigma_0^*/\partial c_0<0$, $\partial\sigma_0^*/\partial(\rho\omega_1)>0$ and $\partial\sigma_0^*/\partial\pi<0$. The effect of $c_0$ on the equilibrium belief, $\partial e^*/\partial c_0=1/(\psi b)$, grows without bound as $\psi\to0$.
    \item In the corner regimes, the premium is decreasing in $\psi$ whenever $C^{\rm def}<\bar\theta C_1+(1-\bar\theta)C_0$ (which implies the corresponding condition at $e(1)$). In Complete Adverse Selection, $\partial\Delta^*/\partial\pi=\psi\,b\,\partial e(1)/\partial\pi<0$: the belief channel of Proposition~\ref{prop:comparative_statics}(i), scaled by $\psi$.
\end{enumerate}
\end{proposition}

\noindent\textit{Proof.} See Appendix~\ref{app:attention}. $\Box$

The comparative statics separate three channels. The \emph{revenue channel} ($\rho\omega_1$) works as in Section~\ref{sec:equilibrium}: a larger effective attention to displayed ratings raises the gain from concealment and the hiding rate. The \emph{belief channel} ($\pi$) keeps its sign but is scaled by $\psi$, because a higher share of informed sellers makes silence more informative only to buyers who think about what silence means. The \emph{attention channel} ($\psi$, $C^{\rm def}$) is new: more attention to silence, or a less optimistic default, lowers the gain from hiding and the hiding rate. These are the margins on which a salience policy acts (Section~\ref{ss:att_policy}). The last statement in part~(i) is the formal reason the single-cost model becomes fragile in apartments: when silence is barely read, a small change in the hiding cost produces a large swing in the hiding rate, because nothing on the buyer side pushes back.

\paragraph{High-quality hiding and non-strategic silence.}
Proposition~\ref{prop:att_incentives}(ii) and Corollary~\ref{cor:att_hq} say when high-quality sellers may hide: only if silence sells above a verified efficient rating. The registry shows similar hiding rates for the two grades, \attSigmaOneApt{} against \attSigmaZeroApt{} percent in apartments and \attSigmaOneHouse{} against \attSigmaZeroHouse{} percent in houses. In apartments silence pays a high-quality seller \attExcessApt{} log points more than her rating, yet she displays, because in Partial Pooling her slack is $b+c_1$ (Proposition~\ref{prop:att_seller_types}). The pattern in the registry is not evidence for strategic hiding by high-quality sellers: house sellers of high quality stay silent although silence costs them about \attHighLossHouse{} log points relative to their rating, so some silence is not strategic at all, for instance listings whose certificate is still pending. The same non-strategic silence is in the hiding rate of low-quality sellers, so the interior hiding rates we use may overstate strategic hiding.

\paragraph{Welfare.}
The welfare accounting of Section~\ref{sec:equilibrium} acquires a third component. As before, the hiding cost $c_0$ is a deadweight loss to the extent it destroys real resources and a transfer to the extent it reflects expected sanctions, informed hiders net zero in Partial Pooling and $\Delta^*-c_0>0$ in Complete Adverse Selection (Proposition~\ref{prop:att_seller_types}), and uninformed sellers earn the non-disclosure premium as a rent. The premium itself now contains a behavioural transfer. Relative to what a fully attentive market would pay, a silent listing sells at a mark-up of
\begin{equation}
    p_{\rm sil}-\big(\bar V-(1-\rho)\,C^{\rm ins}-\rho\,\bar C^{\rm sil}\big)\;=\;\rho\,(1-\omega_0)\,\big(\bar C^{\rm sil}-C^{\rm def}\big),
    \label{eq:att_markup}
\end{equation}
paid by inattentive EPC-sensitive buyers to every silent seller, uninformed ones included. Unlike the allocative distortion of Section~\ref{sec:equilibrium}, which arises because concealed information is not transmitted, this transfer arises because transmitted absence is not read. The model identifies its level only under additional assumptions (Section~\ref{ss:att_identification}); with displayed ratings fully attended and no pre-reform energy pricing it equals $\hat\lambda$, about \attLamApt{} log points per silent apartment sale, and if displayed ratings are not fully attended $\hat\lambda$ is a lower bound.

\subsubsection{The Reform and Identification}
\label{ss:att_identification}

The estimands of Section~\ref{sec:identification} are reform-induced changes, so attention enters through how the reform moved it. For any object $x$, write $\delta x\equiv x^{\rm post}-x^{\rm pre}$ for its change across the July 2021 reform. In this notation $\hat\beta_{1,m}$ estimates $\delta b_m$ and $\hat\Delta_m$ estimates $\delta\Delta_m$. Let the reform raise effective attention to rated listings by $\delta(\rho\omega_1)$ and to silent listings by $\delta(\rho\omega_0)$, and define the \emph{spillover} $\psi_\delta\equiv\delta(\rho\omega_0)/\delta(\rho\omega_1)$, the share of the reform's new attention to energy that reached listings showing no rating. Under the maintained assumption of Section~\ref{sec:identification} that buyers did not price energy before July 2021 (Assumption~\ref{assum:no_pre_pricing}), $\psi_\delta$ equals post-reform relative attention. Normalising the default to zero, equations~\eqref{eq:b_att}--\eqref{eq:premium_att} in differences give
\begin{equation}
    \hat\lambda_m\;\equiv\;\hat\Delta_m-\hat\beta_{1,m}\,\hat e^*_m\;=\;(1-\psi_\delta)\,\hat\beta_{1,m}\,\big(\eta_m-\hat e^*_m\big).
    \label{eq:att_lambda}
\end{equation}

\begin{proposition}[Identification under Limited Attention]
\label{prop:att_identification}
Under Assumptions~\ref{assum:attention} and~\ref{assum:no_pre_pricing} and the maintained assumptions of Section~\ref{sec:identification}:
\begin{enumerate}[label=(\roman*)]
    \item $\hat\rho_m=\hat\beta_{1,m}/\hat\beta^{\rm anchor}_{1,m}$ identifies $\delta(\rho\omega_1)$, the reform-induced increase in effective attention to displayed ratings; with $\omega_1=1$ it is the $\rho_m$ of Section~\ref{sec:identification};
    \item $\hat\Delta_m/\hat\beta_{1,m}$ identifies the perceived cost of silence $C^*_\delta\equiv C^{\rm def}+\psi_\delta(\bar C^{\rm sil}-C^{\rm def})$, in units of certified consumption;
    \item given $C^{\rm def}=0$,~\eqref{eq:att_lambda} point-identifies $\psi_\delta$; without that normalisation, $\psi_\delta\in[0,\,C^*_\delta/\bar C^{\rm sil}]$ and $C^{\rm def}\in[0,C^*_\delta]$;
    \item in Partial Pooling, the hiding cost of the informed low-quality seller equals the post-reform non-disclosure premium, $\delta\Delta$, which $\hat\Delta_m$ estimates.
\end{enumerate}
\end{proposition}

\noindent\textit{Proof.} See Appendix~\ref{app:attention}. $\Box$

The over-identifying wedge is the part of the reform's price effect that silence escaped. Had silent listings received the same increase in attention as rated ones, their price relative to E-to-G listings would have moved by $\hat\beta_1\hat e^*$, as Section~\ref{sec:equilibrium} predicts; they received less, and the shortfall, priced at the cost of the silent pool, is $\hat\lambda$. A spillover of one means the reform made energy something buyers ask about whether or not a rating is shown; a spillover near zero means it made the rating something buyers read but not something they ask for. Part~(iii) states what the design cannot do. The two gaps pin down one number per market, the perceived cost of silence, and low attention with a modest default fits them as well as slightly higher attention with an even more optimistic one. Normalising the default to zero delivers the largest attention compatible with the data.\footnote{The same algebra in levels identifies relative attention at a point in time rather than the spillover. We do not estimate it, because across apartment markets the level gaps carry a large market-specific component that is unrelated to energy and that the differences-in-differences remove: the ratio of mean level gaps and the cross-market slope of one level gap on the other imply very different values, $\attPsiLevROMApt$ and $\attPsiLevSlopeApt$ post-reform. We use levels only for signs.}

Three testable implications follow. By~\eqref{eq:etilde}, $\Delta/b\le\eta$ whenever $C^{\rm def}\ge0$, and since $\eta_m$ comes from the registry alone this is a market-by-market bound on the two estimated gaps. Equation~\eqref{eq:att_lambda} has no constant: a market where the reform did not move the energy-efficiency premium should show no wedge. And the registry belief $\hat e^*_m$ enters~\eqref{eq:att_lambda} with a weight proportional to $\psi_\delta$ (Lemma~\ref{lem:etilde_monotone}), so where attention to silence is low, mismeasurement of the silent pool's composition has little bearing on the estimate.\footnote{At the apartment estimates an error of 0.1 in $\hat e^*$ moves $\hat\psi_\delta$ by about $0.1\times\attSensEApt$. In house markets the sensitivity is about \attSensRatio{} times larger, as in Section~\ref{sec:equilibrium}, and the measurement of the silent pool matters correspondingly more there. For the same reason, the restriction $\sigma_1=0$ imposed in $\hat e^*$ (Section~\ref{sec:beliefs}) matters little for the apartment estimate.}

\subsubsection{Estimates}
\label{ss:att_estimates}

Table~\ref{tab:attention} compares what the core model and the limited-attention model predict for the non-disclosure premium, and reports the attention estimates. We estimate $1-\psi_\delta$ as the ratio of the cross-market sums of the two sides of~\eqref{eq:att_lambda}. Absent selection, and unlike a regression through the origin, this estimator is free of the attenuation and correlated-noise biases that arise when $\hat\beta_{1,m}$ is noisy and shares the E-to-G index with $\hat\Delta_m$. The estimation sample is, however, selected on $\hat\beta_{1,m}>0$, and selecting on a noisy $\hat\beta_{1,m}$ biases both estimators toward full attention: a Monte Carlo calibrated to our markets puts the ratio of sums between \attMCromSelLo{} and \attMCromSelHi{} for a true value of \attMCtrue{} (Appendix~\ref{app:att_estimator}). The apartment estimate is therefore, if anything, an upper bound on attention to silence, and the house estimate may overstate it. We estimate one $\psi_\delta$ per segment rather than per market. A market-specific value is identified by the market's own ratio $\hat\Delta_m/\hat\beta_{1,m}$, but it would fit each premium exactly, place every market in Partial Pooling by construction, and leave no restriction to test, and the market ratios are too noisy to be informative on their own.\footnote{The estimators rank by how much weight they place on markets with a large energy-efficiency premium: the median of market ratios gives $\attPsiMedApt$, the ratio of sums $\attPsiApt$, the regression through the origin $\attPsiOLSApt$, and a regression with a constant, which the model excludes, $\attPsiConstApt$ (intercept $\attIcptApt$, s.e.\ $\attIcptSeApt$).}

\begin{table}[htbp]
\centering
\caption{The Non-Disclosure Premium under the Core Model and under Limited Attention}
\label{tab:attention}
\small\renewcommand{\arraystretch}{1.3}
\begin{tabular*}{\textwidth}{@{\extracolsep{\fill}}lcc@{}}
\toprule\toprule
 & \textbf{Apartments} ($N=45$) & \textbf{Houses} ($N=97$) \\
\midrule
\multicolumn{3}{l}{\textit{A. What the data show (log points)}} \\
Energy-efficiency premium $\hat\beta_1$ & 4.1 & 5.8 \\
Non-disclosure premium $\hat\Delta$ & 8.5 & 3.9 \\
Ratio $\hat\Delta/\hat\beta_1$ (at most one in the core model) & 2.10 & 0.67 \\
Markets above the core-model ceiling $\hat\Delta=\hat\beta_1$ & 37 of 45 & 33 of 97 \\
\midrule
\multicolumn{3}{l}{\textit{B. Predicted non-disclosure premium (log points)}} \\
Core model (full attention): $\hat\beta_1\hat e^*$ & 2.4 & 3.5 \\
\quad Observed minus predicted ($t$) & 6.1 (6.8) & 0.4 (1.1) \\
\cmidrule[1pt]{1-3}
$\triangleright$ \textbf{Limited attention: $\hat\beta_1[\hat\psi\hat e^*+(1-\hat\psi)\eta]$} & \textbf{8.5} & \textbf{3.9} \\
 \quad \textbf{Observed minus predicted} & \textbf{0.0} & \textbf{0.0} \\
\cmidrule[1pt]{1-3}
\multicolumn{3}{l}{\textit{C. Attention to silence}} \\
$\hat\psi_\delta$, ratio of sums [95\% CI] & 0.13 [\ensuremath{-}0.12, 0.37] & 0.95 [0.86, 1.04] \\
$\hat\psi_\delta$, regression through the origin (s.e.) & 0.21 (0.11) & 0.96 (0.04) \\
$t$-statistic, full attention ($\psi_\delta=1$) & 6.9 & 1.1 \\
Identified interval for $\psi_\delta$ ($C^{\rm def}\ge0$) & $[0,\,0.13]$ & $[0,\,0.95]$ \\
Perceived cost of silence $C^*_\delta$ (kWh$_{\rm ep}$/m$^2$/yr) [95\% CI] & 31 [\ensuremath{-}31, 89] & 233 [210, 256] \\
Rational cost of the silent pool $\bar C^{\rm sil}$ (kWh$_{\rm ep}$/m$^2$/yr) & 244 & 245 \\
Energy-cost ratio $\eta$ (weighted) & 2.32 & 2.11 \\
Mean slack in the attention ceiling $\eta\hat\beta_1-\hat\Delta$ ($t$) & 0.90 (1.0) & 8.37 (13.0) \\
Markets above their own attention ceiling & 21 of 45 & 9 of 97 \\
\bottomrule\bottomrule
\end{tabular*}


\vspace{0.5em}
\raggedright
\footnotesize
\textit{Notes:} Difference-in-differences estimates on the estimation sample of Table~\ref{tab:did_results}, log points unless stated; means across markets. Panel~B is the comparison the table is built around. The core model predicts a non-disclosure premium of $\hat\beta_1\hat e^*$, which falls short of the observed premium by \attLamApt{} log points in apartments and by \attLamHouse{} in houses. With limited attention the predicted premium is $\hat\beta_1[\hat\psi\hat e^*+(1-\hat\psi)\eta]$, which matches the observed premium in both segments. In apartments the match uses one parameter, $\hat\psi$, estimated from the ratio of the same two aggregates, so it shows that a single attention parameter is enough and is not an out-of-sample test. The evidence for limited attention is in panel~C: attention to silence is a small fraction of attention to ratings in apartments (\attPsiApt{}) and close to full in houses (\attPsiHouse{}), and full attention is rejected for apartments and not for houses. Aggregates weight markets by $\hat\beta_{1,m}$, the weights implied by the ratio-of-sums estimator. $\eta_m$ and $\bar C^{\rm sil}$ value registry letters at DPE band midpoints; values are provisional pending certified consumption. Confidence intervals from a bootstrap over markets (\attNBoot{} draws); percentile intervals are not truncated to the admissible range $[0,1]$ of $\psi_\delta$. The identified interval allows any default $C^{\rm def}\in[0,C^*_\delta]$. The attention ceiling is $\hat\Delta=\eta_m\hat\beta_{1,m}$ (Section~\ref{ss:att_estimates}).
\end{table}


Figure~\ref{fig:att_contrast} plots the quantity that does not depend on the normalisation of the default, the perceived cost of silence, together with the estimates of $\psi_\delta$.

\begin{figure}[htbp]
    \centering
    \caption{Perceived cost of silence and relative attention, apartments and houses}
    \label{fig:att_contrast}
    \includegraphics[width=\linewidth]{fig/visuals/fig_att_F3_contrast.pdf}
    \begin{minipage}{0.95\linewidth}
        \vspace{0.5em}
        \footnotesize\textit{Notes.} Left: perceived cost of silence $C^*_\delta$ in kWh$_{\rm ep}$/m$^2$/yr with 95\% bootstrap intervals (apartments \attCstarApt{} [\attCstarLoApt, \attCstarHiApt]; houses \attCstarHouse{} [\attCstarLoHouse, \attCstarHiHouse]), against the rational cost of the silent pool (\attCsilApt{} and \attCsilHouse{}). Right: estimates of relative attention $\psi_\delta$ by estimator, and the ratio of sums with the G-band energy midpoint set to 420 or 600 kWh$_{\rm ep}$/m$^2$/yr. The shaded band is the admissible range $[0,1]$. Intervals are not truncated. Apartment intervals are wide, and selection of markets on $\hat\beta_{1,m}>0$ biases the estimates toward full attention.
    \end{minipage}
\end{figure}

% [AUTHOR REWRITE] interpretation of the estimates
Full attention is rejected in apartments and not in houses. The evidence is the contrast between segments, which does not depend on the normalisation of the default: the perceived cost of silence, identified without it, is \attCstarApt{} kWh per square metre and year in apartments (95\% CI \attCstarLoApt{} to \attCstarHiApt{}) against \attCstarHouse{} (\attCstarLoHouse{} to \attCstarHiHouse{}) in houses ($t=\attDiffCstarT$ for the difference). In apartments it is a small fraction of the \attCsilApt{} kWh a fully attentive buyer would assign to the silent pool, so silent apartments are priced roughly as though energy were not part of the transaction. In houses it is close to the rational cost of the pool, which is why the model of Section~\ref{sec:equilibrium} fits that segment. One attention parameter per segment thus accounts for both the rejection in apartments and the fit in houses, without any change to the seller side; its level, unlike the contrast, depends on the normalisation of the default. By Proposition~\ref{prop:att_identification}(iii), the apartment estimate is best read as an upper bound. House prices are consistent with silence being priced at the rational cost of the pool, which is all Section~\ref{sec:equilibrium} requires; they do not show that house buyers are attentive in any absolute sense.

% [AUTHOR REWRITE] fit claim
The extension also changes how well the model predicts the apartment premium, though less than the aggregate suggests (Figure~\ref{fig:att_fit}). It removes the level bias in the predicted premia, the mean residual falling from \attResidCoreApt{} to \attResidExtApt{} log points, but that zero holds by construction, because $\hat\psi_\delta$ is estimated from the same ratio. The root mean squared error falls from \attRmseCoreApt{} to \attRmseExtApt{}, yet net of the level the residual dispersion is unchanged (\attSdResidCoreApt{} against \attSdResidExtApt{}) and the correlation between predicted and observed premia moves only from \attCorrCoreApt{} to \attCorrExtApt{}. The extension therefore explains why silent apartments sell above verified efficient ones, not which markets carry the largest premia. In houses the two models coincide. The evidence for the extension rests on the ceilings of Figures~\ref{fig:att_belief} and~\ref{fig:att_ceilings} (the latter in Appendix~\ref{app:attention}) and the contrast of Figure~\ref{fig:att_contrast}, not on the aggregate fit.

\begin{figure}[htbp]
    \centering
    \caption{Observed and predicted non-disclosure premium, core model and extension}
    \label{fig:att_fit}
    \includegraphics[width=\linewidth]{fig/visuals/fig_att_F1_fit.pdf}
    \begin{minipage}{0.95\linewidth}
        \vspace{0.5em}
        \footnotesize\textit{Notes.} Each market appears twice at the same observed premium $\hat\Delta_m$ (vertical axis, log points): at the core-model prediction $\hat\beta_{1,m}\hat e^*_m$ (hollow) and at the limited-attention prediction $\hat\beta_{1,m}[\hat\psi_\delta\hat e^*_m+(1-\hat\psi_\delta)\eta_m]$ (filled), joined by a thin line. The dashed line marks observed equal to predicted. The box reports the mean residual, the root mean squared error and the correlation of each prediction with $\hat\Delta_m$. The zero mean residual of the extension holds by construction, because $\hat\psi_\delta$ is estimated from the same aggregate ratio. The apartment sample has \attNMarketsApt{} markets and the house sample \attNMarketsHouse{}.
    \end{minipage}
\end{figure}

The attention reading also implies a different concealment cost for apartments: the hiding cost of the marginal hider is the non-disclosure premium itself, \attDthreeApt{} log points, about \attCzeroRatio{} times the \attCzeroBaseApt{} implied by Section~\ref{sec:identification}. A cost of that size is well above the transaction-cost benchmarks of Section~\ref{sec:results}, one reason we keep the baseline estimate as the paper's measure of concealment costs and report this as the alternative reading the extension implies; in houses the two coincide closely (\attCzeroBaseHouse{} and \attDthreeHouse{}).

The model does not take a stance on why apartment buyers fail to attend to energy when no rating is shown. It may be that energy is bundled into service charges. It may be that apartment buyers, often moving from rented housing in which heating was included, have no habit of pricing it. Or it may be that decisions in fast urban markets leave little room to ask. Irrespective of which applies, the implication for prices is the same, and so is the estimate.

Three qualifications bound these results. First, apartment markets cannot be distinguished from the bound of~\eqref{eq:etilde} on average (mean slack \attSlackApt{} log points, $t=\attSlackTApt$): the extension fits them only with a default close to ``no energy cost''. Second, \attNViolApt{} of the \attNMarketsApt{} apartment markets exceed the bound individually; violations are more frequent where $\hat\beta_{1,m}$ is small and the ratio is noisy (\attShareViolBottomApt{} percent of markets in the bottom tercile of $\hat\beta_{1,m}$), but \attShareViolTopApt{} percent of markets in the top tercile also exceed it, so limited attention is the leading force on average but probably not the only one in every market. Third, the level of $\hat\psi_\delta$ depends on the estimator, while the contrast between segments does not.\footnote{The bound depends on how certified consumption is valued. Our $\eta_m$ values registry letters at band midpoints and the open-ended G band at 500 kWh. Valuing G at 420 or 600 kWh moves the apartment estimate of $\psi_\delta$ to $\attPsiGfourtwentyApt$ or $\attPsiGsixhundredApt$, and the number of apartment markets violating the bound to \attNViolGfourtwentyApt{} or \attNViolGsixhundredApt. % [AUTHOR] replace with certified consumption (gate G3) and delete this footnote.
}

Splitting the post-reform window by year, $1-\hat\psi_\delta$ equals \attOneMinusPsiYoneApt, \attOneMinusPsiYtwoApt{} and \attOneMinusPsiYthreeApt{} in apartments and \attOneMinusPsiYoneHouse, \attOneMinusPsiYtwoHouse{} and \attOneMinusPsiYthreeHouse{} in houses. These are imprecise point estimates, computed with the pooled 2022--2024 registry moments; we do not interpret their time pattern without year-specific moments. % [AUTHOR] revisit after G2/G3.

Section~\ref{sec:policy} builds the policy extension on this model.
